\section{MineDojo：把 Minecraft 变成开放式 Agent Research Platform}

有了 recipe，课程先用 MineDojo 把抽象要求落到具体基础设施。它同时提供 environment API、任务生态和 internet knowledge database，使研究者不仅能训练一个 policy，还能研究任务生成、reward learning、知识检索与长期技能。理解这一节时要分清三层：环境能模拟什么、任务怎样定义、互联网材料能提供何种监督。

\subsection{环境、智能体与知识库三件套}

MineDojo 的系统图把开放世界拆成三个互相约束的对象：Minecraft environment 产生视觉与状态，generalist agent 选择 action，internet knowledge base 提供人类示范与文本先验。它不是把网页塞给 agent 就结束，而是要建立 goal、observation、knowledge 和 control 之间的接口。

\lecturefigure{slide-08-minedojo-three-components.jpg}{MineDojo 的开放环境、智能体与互联网知识库}{官方录像恢复幻灯片 V008；00:10:06}

\begin{knowledgebox}{读图：沿数据流而不是沿图标阅读}
从环境出发，智能体收到第一人称画面、inventory 与状态；从知识库出发，它得到视频、Wiki、Reddit 等人类材料；agent 的 action 又改变环境。关键问题是两类输入怎样对齐到同一目标，而不是三块组件是否“都很大”。
\end{knowledgebox}

\lecturefigure{slide-09-three-thousand-tasks-title.jpg}{MineDojo 以 3,000+ 任务组织开放世界研究}{官方录像恢复幻灯片 V009；00:10:24}

“3,000+ tasks”是任务入口，不应读成 generality 已由数量证明。任务规模的价值在于覆盖不同资源、地点、对象、时间跨度和成功条件，让研究者观察 transfer 与 compositionality；若所有任务仍共享同一种浅层模式，数量本身不会产生开放世界能力。

\subsection{Programmatic Tasks 与 Creative Tasks}

MineDojo 的任务空间包含两类互补目标。Programmatic task 可以从 simulator state 自动判断，例如 inventory 是否包含某物、某生物是否被击败；creative task 则可能要求“建造一个好看的房子”或“用有创意的方法灭火”，成功标准含有语义与审美，无法简单写成一个状态谓词。

\lecturefigure{slide-10-task-categories.jpg}{MineDojo 的 simulator、程序化任务与创意任务}{官方录像恢复幻灯片 V010；00:10:48}

\begin{knowledgebox}{术语消化：三类基础设施的职责}
\begin{tabularx}{\textwidth}{@{}L{0.20\textwidth}X X@{}}
\toprule
组件 & 解决的问题 & 与课程主线的关系 \\
\midrule
Versatile simulator & 可观测、可控制、可并行的世界 & 提供主动闭环和可重复实验 \\
Programmatic task & 状态可判定的明确目标 & 适合自动 reward 与大规模训练 \\
Creative task & 难以手写 success predicate 的开放目标 & 迫使系统学习语义 reward 与人类偏好 \\
\bottomrule
\end{tabularx}
\end{knowledgebox}

\lecturefigure{slide-11-inventory-tech-tree-combat.jpg}{Inventory、technology tree 与 combat 形成组合任务空间}{官方录像恢复幻灯片 V011；00:11:18}

这页应先读依赖关系：获得高级物品常要求先采集原料、制作工具、解锁配方并处理危险，因此任务不是独立 label，而是共享一棵技术树。图中的物品、工具和生物只是例子；真正支持迁移的是旧 skill 能被组合成更长 horizon 的新 plan。

\lecturefigure{slide-12-creative-task-examples.jpg}{创意任务缺少简单的自动成功判据}{官方录像恢复幻灯片 V012；00:11:57}

\begin{warningbox}{Creative 不等于“reward 随便给”}
创意任务更难自动评价，但仍需要明确 evidence：目标文本、视频语义、人类偏好、环境约束或 learned reward model。若只用一个模糊分数，policy 可能找到与人类意图无关的 shortcut。开放式评价必须同时报告 reward 来源、失败案例和人工检查。
\end{warningbox}

\lecturefigure{slide-13-firefight-open-ended-demo.jpg}{开放式灭火任务展示语义目标与行动路径}{官方录像恢复幻灯片 V013；00:12:57}

读 demo 时先区分“画面看起来合理”和“任务可验证完成”。视频展示 agent 面对火源、环境结构和工具选择，但单次成功不能给出成功率、泛化范围或 reward robustness。它主要说明开放目标需要把视觉语义、inventory 和多步行动连接起来，而不是只预测下一帧。

\teachervoice{讲者特别区分 programmatic 与 creative tasks：前者容易写 success criterion，后者可能连“什么算完成”都需要语义判断。这个区分直接解释了为什么后面要引入 MineCLIP；learned reward 不是附加模块，而是开放目标进入 RL 的接口。}

\subsection{互联网数据库：从物品表到人类问题}

环境提供体验，但从零 trial-and-error 学会所有配方和策略代价极高。MineDojo 因此把 YouTube、Wiki 与 Reddit 视作人类经验库：视频提供行为与视觉，Wiki 提供结构化事实，论坛提供故障、策略和长尾问题。三种来源的噪声与 grounding 难度不同，不能简单拼接为同质 token。

\lecturefigure{slide-14-internet-knowledge-sources.jpg}{YouTube、Wiki 与 Reddit 构成互补知识来源}{官方录像恢复幻灯片 V014；00:13:06}

\begin{knowledgebox}{读图：三类来源各自补什么}
YouTube 强在动态演示，却没有可靠 action label；Wiki 强在物品、配方和规则，却较少呈现失败恢复；Reddit 强在问题解决与长尾经验，却含有噪声和主观判断。一个 generalist agent 需要 source-aware retrieval，而不是把所有内容都当成同一种监督。
\end{knowledgebox}

\lecturefigure{slide-15-internet-knowledge-scale.jpg}{大规模互联网材料覆盖 Minecraft 的视觉与语义空间}{官方录像恢复幻灯片 V015；00:13:18}

这页的密集缩略图首先说明 coverage，而不是 clean label。读图时比较对象种类、场景与视角多样性；它支持“互联网包含丰富先验”，却不能证明材料已经与 agent state 对齐，更不能证明视频中的行为都正确或可复现。规模必须与去重、质量、版权和任务相关性一起审计。

\lecturefigure{slide-16-minecraft-item-catalog.jpg}{物品、生物与动作目录提供实体层知识}{官方录像恢复幻灯片 V016；00:14:15}

物品目录帮助模型建立 entity vocabulary：哪些对象可采集、可合成、可装备或具有危险。先看类别边界，再看对象与 action 的可供性；知道名称并不等于知道当前视角中哪个像素对应对象，也不等于知道动作前置条件，因此仍需要 perception grounding 和 environment feedback。

\lecturefigure{slide-17-crafting-knowledge.jpg}{Crafting knowledge 把目标物品连接到原料与步骤}{官方录像恢复幻灯片 V017；00:14:24}

制作配方把 declarative fact 变成 procedural constraint。读图时从目标物品反向追踪原料、工具和工作台条件；这种依赖可用于 plan decomposition，也可用于 curriculum。图表支持“互联网可提供结构化步骤”，但执行时仍要检查 inventory、地点与失败恢复。

\lecturefigure{slide-18-potion-recipe-graph.jpg}{药水配方图展示长依赖链与组合计划}{官方录像恢复幻灯片 V018；00:14:33}

这类依赖图的难点不在记住单条 recipe，而在跨多个中间物品维护状态。首先识别目标节点，再沿边回溯所需 ingredient；随后把图上的逻辑顺序映射到世界中的 navigation、collection 与 crafting action。任何缺失资源都可能迫使 agent 重规划。

\lecturefigure{slide-19-wiki-reddit-database.jpg}{Wiki 页面与 Reddit 讨论提供规则和长尾经验}{官方录像恢复幻灯片 V019；00:14:42}

网页与讨论串包含人类解释、替代方案和故障经验，是视频之外的重要 teacher signal。读图时先分事实性页面与意见性帖子，再检查发布时间、版本和上下文。论坛回答可以启发策略，却不应未经验证就执行；environment result 才是当前任务里的最终证据。

\lecturefigure{slide-20-human-open-ended-questions.jpg}{人类提出的开放式问题扩展任务分布}{官方录像恢复幻灯片 V020；00:15:09}

这些问题说明目标分布来自真实玩家，而不是研究者预先设计的有限模板。先看任务是否需要常识、规划或创意，再看能否自动评分。它们支持构建 task proposal 与 retrieval benchmark，但也暴露评价难题：语言目标可能含糊，成功常有多种等价实现。

\teachervoice{讲者把互联网材料称为 reference manual：当人类不知道 Minecraft 中怎样做某件事，也会搜索 Wiki、视频和论坛。Agent 的关键不是“背完整个互联网”，而是在具体 goal 与 state 下找到相关知识，并通过执行结果判断这条建议是否真的适用。}

\lecturefigure{slide-21-generalist-agent-transition.jpg}{从环境与知识库过渡到可学习的 generalist agent}{官方录像恢复幻灯片 V021；00:15:18}

这一过渡页把问题收束为：已有开放环境和人类知识，怎样训练能够理解自然语言目标的 agent？答案首先不是直接预测动作，而是学习一个能把视频与文本目标对齐的 reward model；它把难以手写的语义目标转换成 RL 可使用的反馈。

\subsection{本章小结}

MineDojo 把开放环境、任务生态与互联网知识放进同一平台。Programmatic task 支持自动训练，creative task 暴露语义评价缺口；YouTube、Wiki 与 Reddit 提供互补知识，却必须与当前状态对齐并由执行验证。这为 MineCLIP 的 learned reward 与 Voyager 的 code/memory 路线建立了数据和接口基础。

